% La residencia capitular pertenece a la estructura maestra. Este archivo
% desarrolla su contenido y conserva una segunda etiqueta compatible.
\label{chap:realizacion-espectral-numero-particula}
\index{nombre!réalisation spectrale}
\index{particule!section persistante}
\index{voie!ombre observable}
\index{spectre!réalisation physique}

La fermeture de HMT produit des nombres comme sections généalogiques complètes. La
mécanique dimensionnelle commence lorsqu'une section de cette classe se réalise sur
une voie, transporte un registre, induit un opérateur spectral et acquiert une
carte physique. Ce passage conserve l'antécédent mathématique et change le
codomaine : nombre, état, voie, spectre et particule sont des objets reliés
par des applications explicites et conservent des types différents.

La charnière formelle s'exprime au moyen de
\begin{equation}
 \boldsymbol x
 \xmapsto{\operatorname{sh}}
 \gamma_{\boldsymbol x}
 \xmapsto{\mathcal L}
 \mathcal L_{\boldsymbol x}
 \xmapsto{\Sigma}
 \sigma_{\boldsymbol x}
 \xmapsto{\chi}
 \chi_{\boldsymbol x}
 \xmapsto{\mathcal T}
 \widehat K_{\boldsymbol x}
 \xmapsto{\mathfrak R_{\rm phys}}
 \mathfrak p_{\boldsymbol x}.
 \label{eq:bisagra-numero-espectro-particula}
\end{equation}
Chaque flèche possède un domaine propre et conserve la clé de provenance de la
section initiale.

\section{Correspondance physique du noyau APP--TRIT--TPK}
\label{sec:reaplicacion-fisica-app-trit-tpk}

L'ouverture de la mécanique dimensionnelle parcourt à nouveau le noyau formel avec
un nouveau codomaine. Cette répétition conserve les théorèmes de HMT et déclare
les applications de réalisation physique qui s'ajoutent. Son ordre causal est
\[
 \boxed{\begin{gathered}
 \text{APP comme support de feuillets}
 \longrightarrow \text{TRIT comme orientation temporelle}\\
 \downarrow\\[-1mm]
 \text{cycle unitaire et défaut torsionnel}\\
 \downarrow\\[-1mm]
 \text{coordonnées angulaires duales}
 \longrightarrow \text{particule--voie}
 \end{gathered}}
\]

\subsection{Le tore de Cayley comme support de feuillets physiques}

APP apporte le tore bidimensionnel
\[
 X_9=(\mathbb Z/9\mathbb Z)^2
\]
avec son graphe de Cayley et deux évaluations sur les mêmes sommets. Au
niveau formel, \(\Sigma\) accumule et \(\Pi\) replie. La réalisation physique
attribue à ces évaluations des feuillets et des canaux, en conservant l'application de
provenance qui permet de comparer leurs réponses. Le tore maintient la
périodicité spatiale ; la catégorie des chemins maintient la direction,
la réversibilité et la mémoire du transport.

\subsection{Carte gnomonique euclidienne et atlas des courbures}
\label{sec:carta-gnomonica-equivalencia-ads-ds}

Le gnomon en L complète le noyau multiplicatif \(8\times8\) et sépare
l'intérieur replié de la carte de publication. Dans
la classe nulle, on distingue précisément les deux lecteurs que le quotient
identifie hors de la frontière,
\[
 n\longmapsto n\bmod 9,
 \qquad
 n\longmapsto\operatorname{dr}_9(n),
\]
de sorte que le reste \(0\) et sa publication nonadique \(9\) conservent
des coordonnées distinctes de la même fermeture. Le secteur \(3,6,9\) rend visible
cette transition entre la racine numérique et le reste modulaire.

Dans la réalisation géométrique, le L gnomonique porte la carte euclidienne
\(\mathscr H_0^{(L)}\). Sur celle-ci, les lecteurs inertiel et gravitationnel
descendent à une normalisation commune,
\[
 \frac{\mathcal G_{L}}{\mathcal I_{L}}=1,
\]
qui constitue l'emplacement du principe local d'équivalence. Avant
cette projection, l'état enrichi conserve séparément les lectures
aréale et volumétrique. Le principe d'Ambivalence s'exprime alors par
le rapport interne
\[
 \boxed{
 \frac{\mathcal L_{\rm ar}(x)}{\mathcal L_{\rm vol}(x)}
 =\frac{\pi_{\rm HMT}}{\varphi_{\rm HMT}^{2}}.}
\]
L'égalité euclidienne et le rapport ambivalent appartiennent à des niveaux distincts :
la première est la projection locale sur le L ; le second conserve la
généalogie aréale--volumétrique antérieure à cette projection.

La réalisation métrique s'obtient au moyen d'une famille déclarée d'applications.
Pour
\[
 \mathbb A_s:=\mathbb R[X]/(X^2+s),
 \qquad s\in\{-1,0,+1\},
\]
et une échelle de courbure \(\rho>0\), on fixe
\[
 \mathcal M_{s,\rho}:\mathbb A_s\dashrightarrow
 (M_s,g_{s,\rho}),
 \qquad
 g_{s,\rho}=\mathrm d r^2+S_{s,\rho}(r)^2\mathrm d\theta^2,
\]
avec
\[
 S_{s,\rho}(r)=
 \begin{cases}
  \rho^{-1}\sinh(\rho r),&s=-1,\\
  r,&s=0,\\
  \rho^{-1}\sin(\rho r),&s=+1.
 \end{cases}
 \qquad
 K(M_s,g_{s,\rho})=s\rho^2.
\]
Le domaine radial est \(r\geq0\) pour \(s=-1,0\) ; pour \(s=+1\), la carte
polaire couvre \(0\leq r\leq\pi/\rho\), avec la régularité polaire usuelle à ses
extrémités. La branche \(s=-1\) fournit la section spatiale hyperbolique de la
carte de type anti--de Sitter ; \(s=0\) réalise la carte euclidienne gnomonique ;
et \(s=+1\) fournit la géométrie spatiale positive qui se prolonge dans
les feuilletages de de Sitter. La correspondance complète figure dans les
\cref{sec:ansatz-metrico-tres-curvaturas,sec:atlas-ads-euclideo-ds-hmt}.

La classification algébrique et la famille métrique sont coordonnées au moyen de
la réalisation déclarée \(\mathcal M_{s,\rho}\). Cette flèche constitue
l'ansatz d'interface physique, tandis que
\(K=s\rho^2\) est le calcul exact de la métrique une fois fixés \(s\) et
\(\rho\). La relation quadratique fixe la classe algébrique ; la réalisation
déclarée transporte son signe à la courbure, et \(\rho\) fixe l'échelle
géométrique. Dans la carte négative, la restriction au bord coordonne la
projection archimédienne avec l'incidence exceptionnelle ; dans la carte nulle, les
lecteurs inertiel et gravitationnel descendent à l'équivalence locale ; dans la
carte positive, les feuilletages satisfont \(\ddot a/a=H^2>0\). Cet atlas
relie l'holographie de courbure négative, la physique locale euclidienne et
l'évolution cosmologique accélérée au moyen de domaines et d'applications de changement
déclarés.

\subsection{Représentation temporelle du TRIT : (+1) retour et mémoire, (0) présent de seuil et (-1) ouverture bidirectionnelle}

Une fois déclaré un ordre sur les transitions, le lecteur temporel est
\[
\begin{aligned}
 +1&\longmapsto\text{retour, mémoire retenue et passé},\\
 0&\longmapsto\text{seuil d'évaluation et présent},\\
 -1&\longmapsto\text{ouverture bidirectionnelle et futur}.
\end{aligned}
\]
Le passé désigne la fermeture transportée ; le présent, la frontière
d'évaluation ; le futur, l'arbre des prolongements compatibles. Cette lecture
est la réalisation temporelle des trois algèbres quadratiques déjà construites
dans HMT et conserve l'involution qui échange l'ouverture et le retour.

\subsection{Torsion minimale et dynamique unitaire dans \texorpdfstring{\(\mathbb C^{108}\)}{C108}}

L'invariant dénommé torsion minimale est
\[
 \Delta_4
 =
 \frac{\pi-e\log\pi}{270}
 \left(1+\frac{\pi}{729}\right).
\]
Ses entrées proviennent des sections HMT de \(\pi\) et de \(e\), du canal
\(270\) et de l'intérieur \(729\). Il fonctionne comme mémoire globale du
raffinement avant la spécialisation électronique.

La dynamique temporelle de référence réside dans
\(\mathcal H_{\rm clk}=\mathbb C^{108}\). Avec
\[
 \omega_0=\frac{2\pi}{108\,t_0},
 \qquad
 H_{\rm clk}
 =\hbar_{\rm int}\omega_0
 \sum_{k=0}^{107}k\,|\widetilde k\rangle
 \langle\widetilde k|,
\]
le propagateur satisfait
\[
 U_{\rm step}
 =\exp\!\left(-\frac{\mathrm i}{\hbar_{\rm int}}
 H_{\rm clk}t_0\right),
 \qquad
 U_{\rm step}^{108}=I.
\]
Le cycle unitaire est exact pour les primitives déclarées. La réalisation
comme cristal temporel physique ajoute l'hamiltonien ou le protocole entraîné,
la réponse sous-harmonique primitive, la stabilité sous les perturbations, la
persistance et le mécanisme de protection développés dans le
\cref{chap:regimenes-cuadraticos-tiempo}. Cet emplacement conserve la
hiérarchie et y renvoie la démonstration.

\subsection{Coordonnées angulaires duales et canaux spectraux}

L'incompatibilité somme--produit est transportée vers deux coordonnées angulaires
\(A\) et \(C^*\). Dans la carte archimédienne
\[
 x=\frac{\pi A}{180},
 \qquad
 y=\frac{\pi C^*}{180},
\]
les canaux sont les valeurs propres de l'élément central
\[
 T=e^{-xI-yR}=q_+P_++q_-P_-,
 \qquad
 q_-=e^{-x+y},
 \qquad
 q_+=e^{-x-y}.
\]
La carte est inversible :
\[
 x=-\frac12\log(q_+q_-),
 \qquad
 y=\frac12\log\frac{q_-}{q_+}.
\]
Par conséquent, \(A\) et \(C^*\) sont, respectivement, les coordonnées symétrique et
antisymétrique d'une unique bande spectrale. Le
\cref{chap:canales} contient leur construction et leurs preuves ; la
présente charnière fixe leur position causale avant la particule--voie.

\section{État, voie, forme spectrale et réalisation physique}

\begin{definition}[État fini et section généalogique]
\label{def:estado-finito-seccion-genealogica-md}
Un état fini est un élément
\(x_n\in\mathcal S_n^{\rm surv}\). Une section généalogique est une famille
\(\boldsymbol x=(x_n)_n\) compatible à toute profondeur. L'état est une
restriction de la section ; la section détermine son système complet de
restrictions.
\end{definition}

\begin{definition}[Route observable]
\label{def:ruta-observable-bisagra-md}
Une carte séquentielle de profondeur \(R\) définit l'ombre
\[
 \gamma_{\boldsymbol x}^{[R]}
 :=\operatorname{sh}_R(\boldsymbol x)
 =(c_0\to c_1\to\cdots\to c_R).
\]
La voie conserve la position, le feuillet, l'orientation, la retenue et la mémoire selon la
liste des données de la carte. C'est une présentation ordonnée de la section.
\end{definition}

\begin{definition}[Forme spectrale]
\label{def:forma-espectral-numero-md}
La forme spectrale associée à une section est
\[
 \mathfrak S(\boldsymbol x)
 :=(\gamma_{\boldsymbol x},
    \mathcal L_{\boldsymbol x},
    \sigma_{\boldsymbol x},
    \chi_{\boldsymbol x},
    \widehat K_{\boldsymbol x}),
\]
où \(\mathcal L_{\boldsymbol x}\) est le registre de voie,
\(\sigma_{\boldsymbol x}\) sa signature entière, \(\chi_{\boldsymbol x}\) le
caractère construit sur la signature et \(\widehat K_{\boldsymbol x}\)
l'opérateur obtenu par la hiérarchie des tours. Le spectre
\(\operatorname{Spec}\widehat K_{\boldsymbol x}\) est une sortie interne.
\end{definition}

\begin{definition}[Réalisation physique]
\label{def:realizacion-fisica-forma-espectral}
Une réalisation physique de \(\mathfrak S(\boldsymbol x)\) consiste en
\[
 \mathfrak R_{\rm phys}
 =(\mathcal H_{\rm phys},\rho_{\rm sym},\varsigma_{\bf u},\mathcal O),
\]
où \(\mathcal H_{\rm phys}\) est l'espace des états,
\(\rho_{\rm sym}\) la représentation du groupe de symétrie,
\(\varsigma_{\bf u}\) la carte des unités et \(\mathcal O\) l'ensemble des
observables. La réalisation entrelace les transports de voie avec la dynamique
de \(\mathcal H_{\rm phys}\) et publie des grandeurs avec des unités.
\end{definition}

La forme spectrale est exacte dans le système discret une fois ses
lecteurs fixés. L'identification à une espèce physique appartient à l'application de
réalisation. Cette séparation permet de publier l'opérateur, son spectre, sa
position et son statut, accompagnés de l'opérateur d'entrelacement physique correspondant.

\section{\(4\) grammaires et leurs domaines}

Le passage du nombre à la particule utilise quatre grammaires. Les
trois premières appartiennent au noyau mathématique ; la quatrième appartient à la
représentation statistique choisie.

\begin{center}
\begin{tabular}{@{}lll@{}}
\toprule
\textbf{Grammaire} & \textbf{Opération} & \textbf{Fonction}\\
\midrule
accumulation & feuillet somme d'APP & compose les contributions additives ;\\
repliement & feuillet produit d'APP & compose les échelles et les transports ;\\
orientation & TRIT et conjugaison de voie & distingue le sens et l'anti-section ;\\
occupation & projecteur de la représentation & fixe la multiplicité statistique.\\
\bottomrule
\end{tabular}
\end{center}

L'occupation fermionique se réalise, par exemple, dans une algèbre CAR fixée :
\[
 \{a_i,a_j\}=0,
 \qquad
 \{a_i^\dagger,a_j^\dagger\}=0,
 \qquad
 \{a_i,a_j^\dagger\}=\delta_{ij}I.
\]
De ces prémisses, on déduit
\[
 (a_i^\dagger)^2=0,
 \qquad
 n_i:=a_i^\dagger a_i,
 \qquad
 n_i^2=n_i,
 \qquad
 \operatorname{Spec}(n_i)=\{0,1\}.
\]
Le TPK sélectionne la graduation et transporte la voie ; la représentation CAR
réalise l'exclusion. Cette distinction fixe positivement la provenance
algébrique du principe de Pauli.

\section{Le foncteur de réalisation spectrale}

Soit \(\Omega_{\rm HMT}^{\rm surv}\) la catégorie dont les objets sont des sections
survivantes et dont les morphismes sont des transformations compatibles avec les
restrictions. Soit \(\mathbf{SpecRoute}\) la catégorie des paquets
\((\gamma,\mathcal L,\sigma,\chi,\widehat K)\) et des opérateurs d'entrelacement qui
conservent les clés de voie.

\begin{theorem}[Réalisation spectrale d'une section]
\label{thm:functor-realizacion-espectral-numero}
Les lecteurs de voie, de registre, de signature, de caractère et de tours définissent, dans leur
domaine commun, un foncteur
\[
 \mathfrak S:
 \Omega_{\rm HMT}^{\rm surv}\longrightarrow\mathbf{SpecRoute}.
\]
Pour une restriction compatible de sections, la voie est tronquée, le registre
est restreint et les lecteurs ultérieurs commutent avec cette restriction. Deux
présentations reliées par une symétrie généalogique produisent des formes
spectrales entrelacées.
\end{theorem}

\begin{proof}
L'ombre de voie est définie par lecture séquentielle de l'état et commute
avec la troncature. Le registre concatène des données locales typées et sa
restriction n'élimine que le segment ultérieur. La signature et le caractère
sont calculés au moyen d'homomorphismes déclarés sur ce registre. La
hiérarchie des tours prolonge le caractère par composition. L'identité et la
composition sont préservées à chaque étape.
\end{proof}

Le théorème précise l'expression « le nombre prend une forme spectrale » :
il transporte une section complète vers un paquet de voie et de spectre qui conserve
sa généalogie.

\section{Persistance et définition de la particule}

Pour un préfixe \(s_N\in\mathcal S_N^{\rm surv}\), soit
\[
 \operatorname{Ext}_M(s_N)
 :=\{s_M\in\mathcal S_M^{\rm surv}:\rho_{M,N}(s_M)=s_N\},
 \qquad
 \mathcal L(s_N)
 :=\sup\{M:\operatorname{Ext}_M(s_N)\neq\varnothing\}.
\]

\begin{definition}[Persistance]
\label{def:persistencia-bisagra-numero-particula}
Un préfixe est persistant lorsqu'il appartient à une section globale compatible ;
dans un arbre de prolongements à branchement fini, cette condition équivaut
à \(\mathcal L(s_N)=\infty\).
\end{definition}

\begin{definition}[Particule mathématique]
\label{def:particula-matematica-bisagra}
Soit \(\gamma_{\boldsymbol x}\) une voie persistante et soit
\(\mathcal B_{\gamma_{\boldsymbol x}}\) un fibré discret dont les fibres
portent les degrés internes de la réalisation. Une particule mathématique est
\[
 \mathfrak p=
 (\boldsymbol x,
  \gamma_{\boldsymbol x},
  \mathcal B_{\gamma_{\boldsymbol x}},
  s,n,U,
  \mathfrak S(\boldsymbol x),
  \mathfrak R_{\rm phys}),
\]
où \(s\) est une section compatible du fibré, \(n\) sa fonction
d'occupation, \(U\) le propagateur qui entrelace les transports et
\(\mathfrak R_{\rm phys}\) la réalisation physique déclarée.
\end{definition}

La définition contient trois fermetures :
\begin{enumerate}
\item le nombre est la fermeture globale de préfixes compatibles ;
\item la forme spectrale est la fermeture de la lecture de voie sous le registre,
      la signature, le caractère et les tours ;
\item la particule est la fermeture locale parcourue par une section persistante
      d'un fibré physique.
\end{enumerate}
La relation se résume comme
\[
 \boxed{
 \begin{gathered}
 \text{nombre}=\text{fermeture généalogique globale},
 \\
 \text{particule}=\text{réalisation locale persistante de cette fermeture}.
 \end{gathered}}
\]
L'égalité verbale désigne une dépendance fonctorielle entre types.

\section{Bande orientée et anti-section}

Soit \(\tau=(\tau_0,\ldots,\tau_{11})\in\{-1,0,+1\}^{12}\) le mot
orienté d'une voie et définissons
\begin{equation}
 C_M(\tau):=-\operatorname{rev}(\tau).
 \label{eq:conjugacion-material-bisagra}
\end{equation}
On a \(C_M^2=I\). La bande
\[
 \mathcal B(\tau):=\{\tau,C_M\tau\}
\]
réunit deux orientations d'une même généalogie. Une section réalise la
particule et la section conjuguée réalise l'antiparticule.

\begin{proposition}[Parité de masse et orientation]
\label{prop:paridad-masa-orientacion-bisagra}
Soit \(\mathfrak m(\tau)\) une lecture de masse dépendant d'invariants pairs
de la bande, et soit \(q(\tau)\) une lecture orientée impaire. Alors
\[
 \mathfrak m(C_M\tau)=\mathfrak m(\tau),
 \qquad
 q(C_M\tau)=-q(\tau).
\]
\end{proposition}

\begin{proof}
La première identité exprime l'invariance sous l'involution de bande. La
seconde exprime la parité d'une forme orientée. La relation \(C_M^2=I\)
garantit la compatibilité des deux lectures.
\end{proof}

Particule et antiparticule partagent ainsi le spectre de masse et possèdent
des orientations opposées. La conclusion provient de la parité des lecteurs.

\section{Propagation et somme finie des chemins}

Soit \(U\) un propagateur sur un espace fini d'états. Pour tout
\(R\geq1\), la multiplication matricielle donne
\begin{equation}
 (U^R)_{fi}
 =\sum_{\substack{\gamma:i\to f\\|\gamma|=R}}
   \prod_{\ell\in\gamma}U_\ell.
 \label{eq:feynman-finito-bisagra}
\end{equation}

\begin{theorem}[Développement exact par histoires]
\label{thm:feynman-finito-bisagra}
L'équation \eqref{eq:feynman-finito-bisagra} énumère une fois chaque voie de
longueur \(R\) entre \(i\) et \(f\). Si \(U\) est unitaire, \(U^R\) conserve
la norme et définit une évolution finie unitaire.
\end{theorem}

\begin{proof}
En développant le produit matriciel apparaissent tous les indices intermédiaires
\((x_1,\ldots,x_{R-1})\). Chaque tuple détermine une voie et chaque voie détermine
un tuple. L'unitarité découle de
\((U^R)^\dagger U^R=I\).
\end{proof}

L'amplitude d'onde est la somme cohérente sur la bande des voies ; la
particule est la section persistante dont la propagation réalise cette bande. L'onde
et l'état forment des niveaux complémentaires de la même réalisation.

\section{Horizon fini de prolongement et classe des extensions compatibles de l'état généalogique}

\begin{definition}[Voie virtuelle finie]
\label{def:ruta-virtual-finita-bisagra}
Un préfixe \(s_N\) est virtuel fini s'il existe \(L<\infty\) tel que
\[
 \operatorname{Ext}_L(s_N)\neq\varnothing,
 \qquad
 \operatorname{Ext}_{L+1}(s_N)=\varnothing.
\]
Sa voie transporte le registre, la signature et l'amplitude jusqu'à \(L\), et son horizon de
prolongement se termine à la frontière \(L+1\).
\end{definition}

Dans un arbre à branchement fini, l'existence de prolongements à toute
profondeur équivaut, par le lemme de König, à l'existence d'une section
globale. La persistance et la virtualité sont définies au moyen de la
géométrie des prolongements. La voie virtuelle occupe un état intermédiaire
typé et possède une durée généalogique finie.

\section{Spectre interne, réalisation et espèce}

Pour une multisection \(F\) de voies, l'objet naturel est l'opérateur
\[
 \widehat B_F e_\Gamma
 =\kappa(\Gamma)e_\Gamma,
 \qquad
 \exp((\log R_{\rm act})\widehat B_F).
\]
Son spectre est invariant sous les permutations de la base de la fibre. La carte
physique détermine quelle combinaison de valeurs propres correspond à un observable
scalaire ou matriciel. Trois niveaux doivent être conservés :
\[
 \operatorname{Spec}\widehat B_F,
 \qquad
 \mathfrak R_{\rm phys}(\widehat B_F),
 \qquad
 \text{identification de l'espèce}.
\]
Le premier est une sortie interne ; le deuxième est une réalisation ; le troisième
est une clé dans l'inventaire physique. Un tableau scientifique complet
publie les trois et le morphisme qui les relie.

\section{Foncteur de réalisation physico-spectrale et 1re clôture électronique}

La charnière est fixée par la chaîne
\[
 \boxed{\begin{gathered}
 \text{nombre généalogique}
 \longrightarrow\text{route observable}
 \longrightarrow\text{registre}
 \longrightarrow\text{signature}
 \\
 \text{caractère}
 \longrightarrow\text{tours}
 \longrightarrow\text{opérateur spectral}
 \\
 \text{réalisation physique}
 \longrightarrow\text{particule}
 \end{gathered}}
\]
La première clôture physique spécialise cette chaîne dans la section électronique
sélectionnée par le centre de l'atlas, avec la droite d'action, les lecteurs
bidirectionnels et la carte dimensionnelle en vigueur. Cet emplacement établit le
domaine formel qui permet de comprendre pourquoi l'électron réalise la théorie
du nombre et pourquoi sa généralisation aux autres espèces conserve la voie,
le spectre et la généalogie.

\paragraph{Composition des lecteurs dimensionnels et compatibilité entre leurs invariants mathématiques}
La définition de section persistante et de fibré discret appartient au chapitre
des particules ; la voie, le registre, la signature et le caractère, au chapitre de
l'extension survivante ; la bande, la conjugaison matérielle et la virtualité,
au chapitre des anti-sections ; les opérateurs de fibre et leurs spectres, à la
loi bidirectionnelle des masses. Les quatre grammaires et la fermeture locale du
parcours sont intégrées ici à la structure CAR, à la persistance, au caractère
et à la réalisation physique.
